\documentclass[11pt]{article}
\usepackage[margin=1in]{geometry}
\usepackage[T1]{fontenc}
\usepackage[utf8]{inputenc}
\usepackage{lmodern}
\usepackage{microtype}
\usepackage{amsmath}
\usepackage{amssymb}
\usepackage{booktabs}
\usepackage{enumitem}
\usepackage{hyperref}

\title{\texttt{p\_loop} as the Next Prompt-Profile Objective}
\author{}
\date{2026-03-22}

\begin{document}
\maketitle

\paragraph{Bottom line.}
The next prompt-profile probe should predict direct prompt-level loop probability,
\[
p_{\mathrm{loop}}(x) = \mathbb{E}_r \left[ 1\{\text{rollout } r \text{ loops}\} \right],
\]
from prompt-prefill activations under one fixed model and decode policy. The best
backup objective is \texttt{mean\_relative\_length = E[L / E]}. Keep
\texttt{p\_cap} diagnostic-first, not as the main head.

\paragraph{Why this changed.}
The first real \texttt{GPQA} prompt-profile pilot ruled out the old high-tail
headline target. On the saved \texttt{48}-prompt archive,
\texttt{s\_0.9} collapsed exactly to \texttt{p(cap-hit)}, and
the trained probe was essentially constant-baseline quality. A same-archive
\texttt{mean\_relative\_length} retrain was directionally better, but it still
lost to trivial baselines on MSE/MAE at that pilot scale.

\paragraph{Why \texttt{p\_loop} is the right next scalar.}
The cross-dataset rollout-stat bundle says the failure mode is fundamentally a
terminal-statistics problem:
\begin{itemize}[leftmargin=1.5em]
\item loops are much longer than ordinary wrong rollouts, which are longer than correct ones;
\item looped rollouts are much less accurate;
\item \texttt{P(loop | max-length-hit)} is effectively \texttt{1.0} across all five datasets.
\end{itemize}
That makes \texttt{p\_loop} the cleanest direct prompt-level summary of the
behavior we actually care about: ``under this prompt and this fixed policy, how
often do I enter a looping terminal regime?''

\paragraph{Why not the obvious alternatives.}
\begin{itemize}[leftmargin=1.5em]
\item \textbf{\texttt{majority\_s\_0.5}.} It does show real activation signal on the
  corrected \texttt{GPQA} pilot, but with \texttt{n=4} it is a lossy thresholded
  view of the same rollout counts. Counts \texttt{0}, \texttt{1}, and
  \texttt{2} all map to the same negative label.
\item \textbf{\texttt{s\_t}.} High-threshold tails are too close to cap-hit on the
  current long-cap surface, while lower-threshold tails buy density by mixing in
  more non-loop long answers.
\item \textbf{\texttt{mean\_relative\_length}.} It is the best no-plumbing fallback,
  but it is still a proxy that mixes correct long reasoning, wrong long reasoning,
  and looped long reasoning.
\item \textbf{pure cap risk.} \texttt{p\_cap} is sharp but too narrow. Many loops do
  not hit the cap, especially on \texttt{GPQA} and \texttt{MMLU-Pro}.
\end{itemize}

\paragraph{The key GPQA leakage check.}
On the saved \texttt{GPQA} prompt-profile archive, \texttt{p\_loop} is also the
least prompt-length-correlated of the obvious scalar targets:
\begin{center}
\begin{tabular}{lr}
\toprule
Target & Spearman with prompt length \\
\midrule
\texttt{p\_loop} & \texttt{0.297} \\
\texttt{mean\_relative\_length} & \texttt{0.375} \\
\texttt{p\_cap} & \texttt{0.346} \\
\texttt{s\_0.5} & \texttt{0.341} \\
\texttt{s\_0.6} & \texttt{0.373} \\
\texttt{s\_0.9} & \texttt{0.346} \\
\bottomrule
\end{tabular}
\end{center}
So the direct loop-rate target is not only more aligned than the length-style
proxies; on the current \texttt{GPQA} slice it is also less length-leaky.

\paragraph{Concrete next run.}
\begin{enumerate}[leftmargin=1.5em]
\item \textbf{Split.} Prompt-disjoint, in-distribution only, read per dataset rather
  than pooled first.
\item \textbf{Input.} Last prompt-token prefill activations, comparing:
  \begin{itemize}[leftmargin=1.5em]
  \item one selected-layer MLP probe;
  \item one per-layer MLP ensemble with score aggregation by \texttt{mean\_prob}.
  \end{itemize}
\item \textbf{Target.} \texttt{--target-kind probability --profile-target p\_loop}.
\item \textbf{Backup.} \texttt{--target-kind regression --profile-target mean\_relative\_length}.
\item \textbf{Baselines.} Constant, prompt-token-count only, effective-budget only,
  and prompt-token-count plus effective-budget together.
\end{enumerate}

\paragraph{Evaluation bundle.}
For the main \texttt{p\_loop} head, use the probability-target bundle already in
the repo:
\begin{itemize}[leftmargin=1.5em]
\item Brier;
\item MAE;
\item Spearman;
\item top-\texttt{10\%} and top-\texttt{20\%} capture of realized loop mass.
\end{itemize}
Then attach downstream diagnostics on the same ranked prompts:
\begin{itemize}[leftmargin=1.5em]
\item observed \texttt{p\_cap} in the top-risk buckets;
\item observed correctness in the top-risk buckets;
\item comparison against prompt-length-only and effective-budget-only controls.
\end{itemize}

\paragraph{Implementation note.}
This does not require a new data pipeline. The prompt-profile aggregator already
computes \texttt{p\_loop} for every prompt. The useful code change is just to
expose that existing field as a trainable probability target in the dataset
builder so the next run can call
\texttt{build\_probe\_dataset.py --target-kind probability --profile-target p\_loop}
directly.

\paragraph{Recommendation.}
Stop spending the single-head prompt-profile path on thresholded surrogates first.
Ship \texttt{p\_loop}. If it still fails to beat the leakage baselines on the
larger prompt-disjoint runs, fall back to \texttt{mean\_relative\_length} rather
than another \texttt{s\_t} variant.

\end{document}
